\documentclass[acmsmall,nonacm,screen,10pt]{acmart}
\usepackage{preamble}

%% ── Document metadata ─────────────────────────────────────

\title{Swift Concurrency as a Capability-Region Type System}
\subtitle{Capability, region, and affine transfer in Swift 6.2}

\author{Yasuhiro Inami (@inamiy)}
\thanks{This paper was produced with the assistance of Claude Opus 4.6 and GPT-5.4.}
\affiliation{%
  \institution{try! Swift Tokyo 2026}
  \country{Japan}
}
\date{Apr 13, 2026}

\authorsaddresses{}

\begin{document}

\begin{abstract}
Swift Concurrency exposes programmers to individually sensible rules whose combined behavior is still difficult to predict from surface syntax alone. The hardest cases mix actor isolation, region-based isolation, \code{@Sendable}, \code{sending}, nonisolated async execution, and dynamic \code{@isolated(any)} function values. This paper presents a rule-level account of Swift~6.2 concurrency based on three ingredients: execution capability, value region, and affine environment update. Capability answers \emph{where} a computation runs. Region answers \emph{where} a non-\code{Sendable} value currently lives. Affine update explains \emph{why} a call may preserve access, refine a binding, or consume a value from the caller-side environment. The resulting account clarifies five recurring phenomena: same-isolation \code{sending} does not necessarily consume, ordinary same-isolation calls bind disconnected values into actor regions, cross-isolation results require an explicit \code{sending} asymmetry, \code{@isolated(any)} behaves as a dynamic isolation interface rather than an ordinary actor annotation, and \code{@nonisolated async} after SE-0461 inherits caller execution without collapsing into \code{@concurrent}~\cite{se0461}. The paper validates this model against a focused compiler probe suite covering conversion behavior, task capture, actor-instance closure inference, dynamic isolation identity, and result-transfer corner cases. The appendix prints the complete rule inventory so the paper serves as both an argument and a reference.
\end{abstract}

\keywords{Swift Concurrency, actor isolation, region-based isolation, \texttt{Sendable}, \texttt{sending}, \texttt{@isolated(any)}, type systems.}

\maketitle
{\small Last Update: April 13, 2026}
\bigskip

\section{Introduction}
\label{sec:introduction}

Swift Concurrency gives programmers several ways to talk about isolation: actor annotations, \code{@Sendable}, \code{sending}, \code{isolated} parameters, \code{\#isolation}, \code{@isolated(any)}, and the caller-actor execution model for nonisolated async functions. Each feature has a clear local rationale in the proposal literature, but the combined system is still difficult to predict from surface syntax alone~\cite{se0313,se0414,se0420,se0430,se0431,se0461}. Experienced Swift programmers often know which keyword is present and still fail to predict whether a closure inherits actor context, whether a call consumes a value, or why a sync-looking function still requires \code{await}.

The central difficulty is that the language surface compresses several semantic questions into a single informal word, \emph{isolation}. In practice, three questions must be answered separately.

\begin{itemize}
  \item Where does this computation execute?
  \item Where does this non-\code{Sendable} value currently live?
  \item After this expression, which bindings remain usable in the caller?
\end{itemize}

The motivating example is a minimal closure-conversion failure.

\codefile{snippets/motivation.swift}

The closure literal is written inside a \code{@MainActor} method, yet the parameter type is \code{@Sendable () -> Void}, so the closure body is checked without inherited actor context. The failure is therefore not an arbitrary compiler quirk. It is the visible consequence of a deeper split: sendability constrains capture, while isolation determines execution capability. Once those roles are separated, the diagnostic becomes predictable.

\begin{propositionbox}{Thesis of the paper}
Swift 6.2 concurrency becomes substantially more predictable when it is described with three coordinated notions: execution capability, value region, and affine environment update. Capability determines where code may run. Region determines where a non-\code{Sendable} value is currently connected. Affine update determines whether a call preserves, refines, or consumes the caller-side binding.
\end{propositionbox}

This decomposition draws on ideas from the programming-language theory literature on region types~\cite{tofte1997region,grossman2002cyclone}, static capabilities~\cite{walker2000capabilities,chargueraud2008capabilities}, their intersection~\cite{walker2001regions}, and capability-based concurrency safety~\cite{milano2022fearless}. It does more than explain one closure diagnostic. It explains why same-isolation \code{sending} can preserve access, why ordinary same-isolation calls bind disconnected values into actor regions, why cross-isolation results require explicit \code{sending}, why \code{@isolated(any)} behaves as a dynamic interface rather than a static actor annotation, and why \code{@nonisolated async} after SE-0461~\cite{se0461} is neither actor-hopping folklore nor merely \code{@concurrent} under another name.

\begin{contributionbox}
\begin{itemize}
  \item The paper presents a compact rule kernel that separates execution capability, value region, and affine environment update, and uses that kernel to explain the highest-friction Swift Concurrency behaviors.
  \item It sharpens the conversion story by distinguishing semantic subtyping, one-step contextual coercion, and explicit closure wrapping, with special attention to \code{@isolated(any)}, \code{@concurrent}, and \code{isolated} actor-parameter branches.
  \item It validates the account against an executable compiler-probe suite and prints the complete rule catalog in the appendix so the paper remains useful both as an argument and as a reference artifact.
\end{itemize}
\end{contributionbox}

The rest of the paper proceeds from intuition to formal core and then back to evidence. \Cref{sec:proposal-context} surveys the proposal cluster that defines modern Swift concurrency. \Cref{sec:irregular} explains why the resulting surface model feels irregular. \Cref{sec:model} introduces the judgment form, declaration boundaries, and the capability-region decomposition. \Cref{sec:conversion} characterizes function conversion as a structured family of relations rather than as one flat partial order. \Cref{sec:kernel} isolates the kernel rules that explain transfer, binding, closure inference, and dynamic isolation. \Cref{sec:case-studies} reconnects those rules to concrete case studies. \Cref{sec:validation} summarizes the validation evidence, \Cref{sec:discussion} discusses limits, and \Cref{sec:conclusion} concludes. The appendices print the auxiliary definitions, relation rules, and the full typing-rule inventory.

\section{Proposal Context and Scope}
\label{sec:proposal-context}

The paper studies the proposal cluster that defines modern Swift concurrency behavior. \code{isolated} parameters and actor-instance control come from SE-0313~\cite{se0313}. Region-based isolation and flow-sensitive transfer come from SE-0414~\cite{se0414}. Method and key-path \code{Sendable} inference come from SE-0418~\cite{se0418}. Caller isolation inheritance and \code{\#isolation} come from SE-0420~\cite{se0420}. Explicit \code{sending} parameters and results come from SE-0430~\cite{se0430}. Dynamic \code{@isolated(any)} function values come from SE-0431~\cite{se0431}. The caller-actor execution model for nonisolated async functions is revised by SE-0461~\cite{se0461}.

The goal is not to restate every rule those proposals contain. The goal is to extract the smallest account that still explains the behaviors programmers most often find irregular:

\begin{itemize}
  \item closure conversion and the loss or preservation of inherited actor context,
  \item the difference between same-isolation binding and cross-isolation transfer for non-\code{Sendable} values,
  \item the boundary between \code{@nonisolated async} and \code{@concurrent},
  \item dynamic isolation through \code{@isolated(any)},
  \item and caller-side evidence mechanisms such as \code{isolated} parameters and \code{\#isolation}.
\end{itemize}

The resulting model is intentionally selective. It is not a full mechanization of Swift, and it does not try to formalize every parser rule, standard-library API, or SIL optimization. It instead aims for explanatory leverage: a vocabulary strong enough to make the recurring diagnostics predictable.

\section{Why the Surface Model Feels Irregular}
\label{sec:irregular}

The most common informal explanation of Swift Concurrency is that code is either isolated or nonisolated. That explanation is too coarse for at least three reasons.

First, function values and function bodies live on different semantic layers.

\codefile{snippets/handler-location.swift}

The property \code{handler} lives in the region of the enclosing declaration, while the function body that the property stores runs according to the isolation annotation on the function type. A single phrase such as ``the closure is on \code{@FooActor}'' is therefore ambiguous. One must distinguish storage location from execution capability.

Second, accessibility and capturability are not the same relation. A task-isolated value may be directly readable in some contexts and still be illegal to capture into an actor-isolated closure. This distinction is crucial for understanding nonisolated async parameters after SE-0461 and for explaining why \code{@isolated(any)} closures can capture some values that \code{@MainActor} closures cannot~\cite{se0461}.

Third, boundary crossing is not the whole story. Same-isolation ordinary calls do not consume non-\code{Sendable} values, but they may refine a disconnected value into an actor-bound region. Same-isolation \code{sending} calls preserve the original region instead of binding it. Cross-isolation calls may consume the argument, while cross-isolation results require a separate explicit \code{sending} guarantee. The model therefore needs more than a boolean boundary test; it needs a caller-side environment that changes over time.

\begin{definitionbox}{Three semantic relations}
The paper distinguishes three relations that are often conflated in surface reasoning. $\Accessible(@\kappa)$ determines which regions can be used directly in a capability~$@\kappa$. $\Capturable(@\kappa)$ determines which regions may be captured into a closure checked under~$@\kappa$. $\canSend(@\kappa, @\iota, [\sending])$ determines when a call site performs an explicit or implicit transfer.
\end{definitionbox}

These distinctions also explain why global actors and actor instances behave differently under closure inference. A global-actor closure can inherit actor context without capturing a concrete actor value. An actor-instance closure keeps instance isolation only if the isolated parameter or \code{self} is actually captured. That capture-sensitive branch is easy to miss in proposal prose, but it is essential for matching compiler behavior.

The next section turns these informal observations into a judgment form and a small set of auxiliary definitions.

\section{Formal Ingredients}
\label{sec:model}

The paper uses two semantic axes and one update discipline.

\begin{itemize}
  \item Capability answers \emph{where} the current computation runs.
  \item Region answers \emph{where} a non-\code{Sendable} value is currently connected.
  \item Affine environment update answers \emph{what} remains usable after evaluation.
\end{itemize}

The key claim is that Swift Concurrency becomes predictable only when these three roles are kept distinct.

\subsection{Judgment and Declaration Boundaries}

\begin{rulebox}{Expression judgment}
\[
  \Gamma;\; @\kappa;\; \alpha \vdash e : T \at \rho \dashv \Gamma'
\]

\begin{itemize}
  \item $\Gamma$: input environment
  \item $\Gamma'$: output environment
  \item $@\kappa$: current execution capability
  \item $\alpha$: ambient mode, either $\mathit{sync}$ or $\mathit{async}$
  \item $\rho$: result region
\end{itemize}

The final component is essential because evaluation may preserve a binding, refine its region, or remove it from the caller-side environment.
\end{rulebox}

Swift function declarations are not expression terms, so the formalization also needs a declaration judgment that says how a body is checked. A declaration contributes two pieces of information to its body: the body capability~$@\kappa$ and the ambient mode~$\alpha$. For a function type annotation $@\varphi = @\sigma\; @\iota$, the body capability is obtained from $\toCapability(\mathrm{proj}_\iota(@\varphi))$, while $\alpha$ is determined by the presence or absence of \code{async}. This is where the paper introduces \code{await}-permission; it is not encoded as a separate effect system.

\begin{definitionbox}{Function annotations and body capability}
\[
\begin{aligned}
  @\sigma &::= \cdot \mid @\textsf{Sendable} \\
  @\iota &::= @\kappa \mid @\textsf{isolated(any)} \mid @\textsf{concurrent} \\
  @\varphi &::= @\sigma\; @\iota
\end{aligned}
\]

\[
\begin{aligned}
  \toCapability(@\kappa) &= @\kappa \\
  \toCapability(@\textsf{isolated(any)}) &= @\textsf{nonisolated} \\
  \toCapability(@\textsf{concurrent}) &= @\textsf{nonisolated}
\end{aligned}
\]
\end{definitionbox}

This split between $@\sigma$ and $@\iota$ matters immediately. \code{@Sendable} is not an execution mode. It constrains capture. By contrast, $@\iota$ determines the body capability. Both \code{@concurrent async} and \code{nonisolated(nonsending) async} therefore type-check their bodies under $@\textsf{nonisolated}$, even though their call-site semantics differ later~\cite{se0461}.

\subsection{Capability, Region, and Two Access Relations}

Capability alone does not explain what remains legal after a call or an assignment. The model also needs a value-side location discipline. The region component is inspired by Tofte and Talpin's region-based memory management~\cite{tofte1997region}, while the capability component follows the static-capability tradition of Walker, Crary, and Morrisett~\cite{walker2000capabilities}. The interplay between regions and linear ownership draws on Walker and Watkins~\cite{walker2001regions}.

\[
\begin{aligned}
  \rho_{\mathit{ns}} &::= \disconnected \mid \isolated(a) \mid \task \mid \invalid \\
  \rho &::= \rho_{\mathit{ns}} \mid \mathunderscore
\end{aligned}
\]

$\disconnected$ means not yet bound to an isolation domain. $\isolated(a)$ means connected to a concrete actor region. $\task$ denotes task-local non-\code{Sendable} state. $\mathunderscore$ is the normalized region tag for \code{Sendable} values. The normal form condition is therefore simple: $T : \textsf{Sendable}$ if and only if $\rho = \mathunderscore$.

Two relations are then needed.

\begin{itemize}
  \item $\Accessible(@\kappa)$ determines which regions may be used directly under capability $@\kappa$.
  \item $\Capturable(@\kappa)$ determines which regions may be closed over by a closure checked under $@\kappa$.
\end{itemize}

This distinction is one of the paper's critical refinements. Actor-isolated code may directly access task-bound values in some contexts, but an actor-isolated closure may not capture a \code{task} value because the captured state could outlive the dynamic situation that made the direct access safe. That is why nonisolated async parameters can be captured by nonisolated and \code{@isolated(any)} closures yet still be rejected by \code{@MainActor} closures after SE-0461~\cite{se0461}.

\subsection{Region Merge and Capability-Region Glue}

\codefile{snippets/region-binding.swift}

The snippet above illustrates why region and capability cannot be merged into one notion. The value \code{x} begins disconnected. Storing it into actor-isolated state does not change its nominal type, but it does change the future facts that the type checker must remember. The environment must therefore record a refinement from $\disconnected$ to an actor-bound region.

\paperfigure[0.78]{figures/region-merge-semilattice.png}{Region merge for non-\code{Sendable} values. A disconnected value can be refined into an actor or task region, while incompatible branches join to $\invalid$, which corresponds to a compile-time failure.}

The operation that connects capability and region is $\toRegion(@\kappa)$. It answers a narrow but important question: if a same-isolation operation binds a value to the current execution context, which region should the value acquire? For a concrete actor capability the answer is that actor's isolated region; for $@\textsf{nonisolated}$ the answer remains $\disconnected$. This definition is what lets the formalization state same-isolation binding without conflating execution and storage.

\begin{propositionbox}{Why the output environment is unavoidable}
Region merge is not a side remark. It is the mechanism that explains why local code changes future legality. The result of an assignment or a same-isolation call is often \code{Sendable} and therefore normalized to region $\mathunderscore$, while the interesting fact survives only in the updated environment. A model without $\Gamma'$ can describe the value just produced, but it cannot describe why a later transfer becomes illegal.
\end{propositionbox}

\subsection{Transfer Is a Predicate, Not a Slogan}

Swift offers both explicit and implicit transfer. Explicit transfer is written with \code{sending}. Implicit transfer occurs when a non-\code{Sendable} value crosses to a distinct isolation domain. The paper packages both cases into one predicate:

\[
  \canSend(@\kappa, @\iota, [\sending])
\]

$\canSend$ holds either because the parameter is explicitly marked \code{sending}, or because the callee isolation $@\iota$ is distinct from the caller capability $@\kappa$ and therefore induces an implicit cross-isolation transfer. This is another place where surface syntax hides a deeper commonality: the important fact is not merely whether the keyword \code{sending} appears, but whether the call site transfers responsibility for the value.

This viewpoint also clarifies the asymmetry between parameters and results. Parameters can be consumed from the caller by shrinking the environment. Results travel in the opposite direction. A cross-isolation non-\code{Sendable} result is therefore legal only if the callee explicitly promises a disconnected \code{sending} result. The next two sections show how these auxiliary definitions govern conversion, calls, closure inference, and dynamic isolation.

\section{Conversion Is Not One Partial Order}
\label{sec:conversion}

Function conversion is where informal explanations of Swift Concurrency most often collapse. It is tempting to ask whether one isolated function type is a subtype of another and stop there. The experimental conversion corpus shows that this is too simple. The formalization needs three distinct notions.

\begin{itemize}
  \item semantic subtyping, which may compose transitively,
  \item one-step contextual coercion, which represents the direct edges the compiler accepts without rebuilding the function value,
  \item and explicit closure wrapping, where a new function is constructed to mediate between source and target.
\end{itemize}

\begin{definitionbox}{Why the distinction matters}
Some successful local conversions are genuine subtype or direct coercion edges. Others succeed only because the programmer writes a fresh closure such as \code{\{ await f(\$0) \}} or because contextual typing synthesizes an equivalent wrapper. Treating these cases as one relation hides the structure that the compiler actually implements.
\end{definitionbox}

The conversion graphs therefore do not describe a single flat lattice. They describe several interacting families: sync conversion, async conversion, dynamic-isolation conversion around \code{@isolated(any)}, and actor-parameter branches that cannot be erased by direct coercion.

\paperfigure[0.92]{figures/func-conversion-sync.png}{Selected sync conversion behavior. The graph contains genuine subtype and coercion edges, but not every useful sync-to-sync change is represented. In particular, actor-parameter branches and deprecated \code{@isolated(any)} erasure sit outside a simple monotone order.}

\paperfigure[0.92]{figures/func-conversion-async.png}{Selected async conversion behavior. Async introduces a richer direct-coercion space: non-\code{Sendable} \code{@nonisolated}, \code{@concurrent}, and \code{@isolated(any)} form one practical cluster, while \code{@Sendable} variants form another. Actor-parameter branches remain a separate sink that usually requires closure wrapping.}

Three observations matter most for the rest of the paper.

\subsection{Sync and Async Have Different Structural Freedom}

In the sync world, unrelated actor-isolated functions are not interchangeable because there is no suspension point that could bridge the actor boundary. Sync conversion therefore stays narrow. By contrast, async function values admit a broader coercion space because an async call can incorporate suspension and actor hopping when required. This is why \code{@MainActor async} behaves far more flexibly than \code{@MainActor sync}, and why the async graph contains genuinely bidirectional clusters that do not exist in sync form~\cite{se0431,se0461}.

\subsection{\code{@concurrent} and \code{@nonisolated async} Are Not Identical}

Both \code{@concurrent async} and \code{@nonisolated async} bodies type-check under $@\textsf{nonisolated}$, but they represent different operational stories. \code{@nonisolated async} after SE-0461 inherits caller execution, whereas \code{@concurrent} is explicitly actor-independent and performs stricter sendability checking at actor-capable call sites~\cite{se0461}. The conversion graphs partly overlap because both inhabit the non-actor-isolated async world, but their call rules diverge later.

\subsection{\code{@isolated(any)} Is a Dynamic Interface, Not a Static Actor Label}

\code{@isolated(any)} does not merely mean ``some actor annotation I do not care about.'' It packages actor identity dynamically. That is why a sync-looking \code{@isolated(any)} value may still require \code{await}, and why actor identity can be preserved or erased depending on the conversion path that produced the value~\cite{se0431}.

\codefile{snippets/isolated-any-await.swift}

The example above is not a corner case. It expresses the central lesson of dynamic isolation: a value of type \code{@isolated(any) () -> Void} carries enough runtime information to require cross-isolation calling conventions, even when no static actor name appears at the call site.

\begin{propositionbox}{The actor-parameter branch is a sink, not a peer}
Function types with an \code{isolated} actor parameter are structurally different from zero-argument actor-isolated functions. They carry a concrete actor witness as part of the function value's calling convention. Direct coercion may forget \code{@Sendable}, but it cannot erase that actor parameter. Reaching or leaving this branch generally requires explicit closure wrapping rather than ordinary subtype reasoning.
\end{propositionbox}

The conversion section therefore sets up the remainder of the paper in two ways. First, it explains why the call rules must distinguish same-isolation, cross-isolation, \code{@concurrent}, and dynamic-isolation cases instead of relying on one generic ``isolated function'' intuition. Second, it explains why validation must check both direct coercion edges and identity-erasing construction paths. The next section isolates the rules that perform that work.

\section{A Kernel of Typing Rules}
\label{sec:kernel}

Appendix~\ref{sec:appendix-full} prints the complete rule inventory. The main body needs only the rules that explain the recurring behaviors. Those rules fall into four groups: access and binding, transfer, closure inference, and dynamic isolation.

\subsection{Access and Same-Isolation Binding}

\begin{rulebox}{Variable access depends on capability and region}
\typerule{var}
  {x : T \at \rho \in \Gamma \\ \rho \in \Accessible(@\kappa)}
  {\Gamma;\; @\kappa;\; \alpha \vdash x : T \at \rho \dashv \Gamma}
\end{rulebox}

The \textsc{var} rule is intentionally mundane. Its purpose is to make the separation explicit: access is legal not merely because $x$ has type~$T$, but because the current capability can directly access the region in which $x$ currently lives.

\begin{rulebox}{Same-isolation ordinary call refines the caller environment}
\typerule{call-same-nonsendable-merge}
  {\Gamma;\; @\kappa;\; \alpha \vdash f : @\sigma\; @\kappa\; (A) \to B \at \rho_f \dashv \Gamma_1 \\
   \Gamma_1;\; @\kappa;\; \alpha \vdash \mathit{arg} : A \at \rho_a \dashv \Gamma_2 \\
   A : {\sim}\textsf{Sendable} \\
   \rho_a \in \Accessible(@\kappa) \\
   \rho'_a = \rho_a \sqcup \toRegion(@\kappa) \\
   \rho_b \in \Accessible(@\kappa),\; (B : \textsf{Sendable} \Rightarrow \rho_b = \mathunderscore)}
  {\Gamma;\; @\kappa;\; \alpha \vdash f(\mathit{arg}) : B \at \rho_b \dashv (\Gamma_2[\mathit{arg} \mapsto A \at \rho'_a])}
\end{rulebox}

This rule explains one of the paper's most important phenomena: same-isolation use does not imply stasis. A disconnected value passed to an ordinary same-isolation non-\code{sending} parameter may become actor-bound after the call. The value is still usable in the caller, but not under the same preconditions as before. That is why later transfer can fail even without any explicit consume. The result region~$\rho_b$ generalizes naturally: when $B$ is \code{Sendable} it normalizes to~$\mathunderscore$; when $B$ is non-\code{Sendable} it may be \code{disconnected} (fresh value) or actor-bound (actor state).

The same join operation also governs assignment and aliasing through \textsc{region-merge}. The paper does not treat assignment as destruction; it treats it as refinement of the environment's region facts.

\subsection{Transfer Depends on More Than the Keyword \code{sending}}

\begin{rulebox}{Same-isolation \code{sending} without consumption}
\typerule{call-nonsendable-noconsume}
  {\Gamma;\; @\kappa;\; \alpha \vdash f : @\sigma\; @\iota\; ([\sending]\; A)\; \alpha' \to B \at \rho_f \dashv \Gamma_1 \\
   \Gamma_1;\; @\kappa;\; \alpha \vdash \mathit{arg} : A \at \rho_a \dashv \Gamma_2 \\
   A : {\sim}\textsf{Sendable},\; B : \textsf{Sendable},\; @\iota \notin \{@\textsf{concurrent}\} \\
   \canSend(@\kappa, @\iota, [\sending]) \\
   \rho_a \in \Accessible(@\kappa) \\
   \isActorIsolated(@\kappa) \land @\kappa = @\iota}
  {\Gamma;\; @\kappa;\; \alpha \vdash [\code{await}]\; f(\mathit{arg}) : B \at \mathunderscore \dashv \Gamma_2}
\end{rulebox}

\begin{rulebox}{Cross-isolation transfer consumes the caller binding}
\typerule{call-nonsendable-consume}
  {\Gamma;\; @\kappa;\; \alpha \vdash f : @\sigma\; @\iota\; ([\sending]\; A)\; \alpha' \to B \at \rho_f \dashv \Gamma_1 \\
   \Gamma_1;\; @\kappa;\; \alpha \vdash \mathit{arg} : A \at \disconnected \dashv \Gamma_2 \\
   A : {\sim}\textsf{Sendable},\; B : \textsf{Sendable},\; @\iota \notin \{@\textsf{concurrent}\} \\
   \canSend(@\kappa, @\iota, [\sending]) \\
   \neg(\isActorIsolated(@\kappa) \land @\kappa = @\iota)}
  {\Gamma;\; @\kappa;\; \alpha \vdash [\code{await}]\; f(\mathit{arg}) : B \at \mathunderscore \dashv (\Gamma_2 \setminus \{\mathit{arg}\})}
\end{rulebox}

These two rules express the paper's central contrast. \code{sending} marks a transfer-capable boundary, not an unconditional consume~\cite{se0430}. If caller and callee share the same concrete actor isolation, the value remains in the caller environment with its original region. If that same-isolation proof fails, the argument must be disconnected and is removed from the caller environment after the call.

\begin{definitionbox}{Parameter and result positions are asymmetric}
For parameters, cross-isolation transfer is safe because the caller can relinquish the binding and the type system can record that fact by shrinking~$\Gamma$. Results move in the opposite direction. A cross-isolation non-\code{Sendable} result is therefore legal only when the callee explicitly promises \code{sending}~$B$, which yields a disconnected value in the caller. Without that promise, derivation fails.
\end{definitionbox}

This result-side asymmetry is not cosmetic. It explains why cross-isolation calls allow implicit transfer of disconnected arguments yet reject ordinary non-\code{Sendable} return values without explicit \code{sending}~\cite{se0430}.

\begin{propositionbox}{\code{@concurrent}, \code{@nonisolated async}, and \code{@nonisolated sync} are separate transfer stories}
\code{@concurrent async} requires a non-\code{Sendable} argument to be disconnected, but it does not consume that argument because the call is modeled as a callability check rather than as actor-to-actor transfer. By contrast, \code{@nonisolated async} after SE-0461 inherits caller execution for non-\code{sending} calls and therefore preserves the caller-side argument region. Its results are more subtle: a freshly constructed non-\code{Sendable} result may be disconnected, but a returned input may remain task-bound rather than transfer-ready~\cite{se0461}. A \code{@nonisolated sync} function called from an actor-isolated context is the simplest case: it executes on the caller's executor with no isolation boundary, requiring neither \code{await} nor argument consumption, so the environment is preserved exactly as-is.
\end{propositionbox}

\subsection{Closure Inference Splits on Sendability and Capture}

\begin{rulebox}{Parent-inheriting closure}
\typerule{closure-inherit-parent}
  {@\kappa_{\mathit{eff}} = \effectiveClosureCapability(@\kappa, \{e\}) \\
   \Gamma_{\mathit{captured}} \subseteq \Gamma \land \neg\isPassedToSendingParameter(\{e\}) \\
   \forall (y : U \at \rho) \in \Gamma_{\mathit{captured}},\; \rho \in \Capturable(@\kappa_{\mathit{eff}}) \\
   \Gamma_{\mathit{captured}};\; @\kappa_{\mathit{eff}};\; \alpha' \vdash e : B \at \rho_{\mathit{ret}} \dashv \Gamma'_{\mathit{cl}} \\
   \rho_{\mathit{closure}} = \bigsqcup\{\rho \mid (y : T \at \rho) \in \Gamma_{\mathit{captured}} \land T : {\sim}\textsf{Sendable}\}}
  {\Gamma;\; @\kappa;\; \alpha \vdash \{e\} : @\kappa_{\mathit{eff}}()\; \alpha' \to B \\
   \at \rho_{\mathit{closure}} \dashv \Gamma}
\end{rulebox}

\begin{rulebox}{\code{@Sendable} closure without parent inheritance}
\typerule{closure-no-inherit-parent}
  {\Gamma_{\mathit{captured}} \subseteq \Gamma
   \quad \neg\inheritActorContext(\{e\}) \land \neg\isPassedToSendingParameter(\{e\}) \\
   \isAllSendable(\Gamma_{\mathit{captured}})
   \quad @\kappa_{\mathit{cl}} = \toCapability(@\iota) \\
   \Gamma_{\mathit{captured}};\; @\kappa_{\mathit{cl}};\; \alpha' \vdash e : B \at \rho_{\mathit{ret}} \dashv \Gamma'_{\mathit{cl}}}
  {\Gamma;\; @\kappa;\; \alpha \vdash \{e\} : @\textsf{Sendable}\; @\iota\; ()\; \alpha' \to B \at \mathunderscore \dashv \Gamma}
\end{rulebox}

This pair of rules explains the motivating example from the introduction. Non-\code{@Sendable} closures inherit effective parent capability. \code{@Sendable} closures do not; they are checked under the capability implied by their own annotation. The split is conceptually simple, but it has two subtle consequences.

\begin{itemize}
  \item \code{@Sendable} constrains capture rather than directly describing execution. That is why a \code{@Sendable @MainActor} closure is legal: it is sendable as a value, yet its body still executes under \code{@MainActor}.
  \item Actor-instance inheritance is capture-sensitive. A global-actor closure keeps global actor isolation without capturing a concrete actor value, but a closure inside an actor method or under an \code{isolated} parameter keeps instance isolation only if \code{self} or the isolated parameter is actually captured.
\end{itemize}

\begin{propositionbox}{Sending a closure revives the same-isolation versus transfer split}
When a closure itself is passed through a \code{sending} parameter, the same distinction appears at the capture level. Same concrete actor inheritance, such as \code{Task \{ \ldots \}} in an actor-isolated caller with \code{@\_inheritActorContext}, preserves captures. Transfer cases, such as \code{Task.detached} or nonisolated callers, require non-\code{Sendable} captures to be disconnected and consume them from the caller environment.
\end{propositionbox}

This is why the folklore claim ``tasks always consume captures'' is false. What matters is whether the sent closure still has same concrete actor isolation evidence.

\subsection{Structured Concurrency: \code{async let}}

\begin{rulebox}{\code{async let} declares a child task with temporary capture consumption}
\typerule{async-let}
  {\Gamma_{\mathit{captured}} \subseteq \Gamma \\
   \forall (y : U \at \rho) \in \Gamma_{\mathit{captured}}.\; (U : \textsf{Sendable}) \lor (U : {\sim}\textsf{Sendable} \land \rho = \disconnected) \\
   \Gamma_{\mathit{captured}};\; @\textsf{nonisolated};\; \mathit{async} \vdash \mathit{expr} : T \at \rho_{\mathit{init}} \dashv \Gamma'_{\mathit{child}} \\
   \Gamma_{\mathit{sent}} = \Gamma \setminus \{y \mid (y : U \at \disconnected) \in \Gamma_{\mathit{captured}} \land U : {\sim}\textsf{Sendable}\} \\
   \rho_x = (\text{if } T : \textsf{Sendable} \text{ then } \mathunderscore \text{ else } \disconnected)}
  {\Gamma;\; @\kappa;\; \mathit{async} \vdash \mathbf{async\; let}\; x = \mathit{expr} \dashv \Gamma_{\mathit{sent}},\; x : T \at \rho_x}
\end{rulebox}

\begin{rulebox}{\code{async let} access restores non-consumed captures}
\typerule{async-let-access}
  {(x : T \at \rho) \in \Gamma \quad \isAsyncLet(x) \\
   \Gamma' = \undosend(\Gamma, x)}
  {\Gamma;\; @\kappa;\; \mathit{async} \vdash \mathbf{await}\; x : T \at \rho \dashv \Gamma'}
\end{rulebox}

These two rules capture the novel semantics of \code{async let}. The initializer is type-checked under $@\textsf{nonisolated}; \mathit{async}$, creating a child task boundary. Because the conclusion already requires ambient mode~$\mathit{async}$, \code{async let} is available only in async contexts. Non-\code{Sendable} captures must be disconnected and are temporarily consumed (removed from~$\Gamma$). Unlike \code{Task.init}, which can inherit the caller's actor isolation through \code{@\_inheritActorContext} and thereby preserve captures, \code{async let} always acts as a nonisolated boundary.

The $\undosend$ function distinguishes \code{async let} from ordinary \code{sending}-based transfer. At \code{await}~$x$, the SIL region analysis restores captures that were not consumed by cross-isolation within the child task body. Formally: $\undosend(\Gamma, x) = \Gamma \cup \{y : U \at \disconnected \mid y \in \mathit{sent}_{\text{async-let}}(x) \land y \in \mathrm{dom}(\Gamma'_{\mathit{child}})\}$. Here $\mathrm{dom}(\Gamma)$ denotes the domain of environment~$\Gamma$, i.e.\ the set of variable names $y$ such that $\exists T, \rho.\; y : T \at \rho \in \Gamma$. When the initializer makes only same-isolation calls (e.g.\ a nonisolated function), $\Gamma'_{\mathit{child}}$ retains all captures and $\undosend$ fully restores them. When the initializer crosses isolation boundaries (e.g.\ calling a \code{@MainActor} function), those captures are consumed within $\Gamma'_{\mathit{child}}$ and remain unavailable after \code{await}.

\subsection{Dynamic Isolation Requires an Explicit Bridge}

\begin{rulebox}{\code{@isolated(any)} sync-to-async bridge}
\typerule{isolated-any-to-async}
  {\Gamma;\; @\kappa;\; \alpha \vdash f : @\sigma\; @\textsf{isolated(any)}\; () \to B \at \rho_f \dashv \Gamma' \\
   \neg\isActorIsolated(@\iota_2) \\
   \rho' = (\text{if } @\sigma = @\textsf{Sendable} \text{ then } \mathunderscore \text{ else } \rho_f)}
  {\Gamma;\; @\kappa;\; \alpha \vdash f : @\sigma\; @\iota_2\; ()\; \mathit{async} \to B \at \rho' \dashv \Gamma'}
\end{rulebox}

The bridge exists because \code{@isolated(any)} values carry actor identity dynamically~\cite{se0431}. A sync-looking value of type \code{@isolated(any) () -> B} must still be callable through cross-isolation conventions, and those conventions are async. The same dynamic information also explains the observable property \code{f.isolation}: some construction paths preserve actor identity, while others erase it permanently by converting through an ordinary nonisolated function type.

Taken together, the rules in this section explain why Swift Concurrency diagnostics often feel related even when they appear in different syntactic corners. Access, binding, transfer, capture, and dynamic calling conventions are not separate ad hoc mechanisms. They are different projections of one capability-region discipline with affine update.

\section{Case Studies}
\label{sec:case-studies}

The rule kernel is easier to trust once it is tied back to concrete examples. This section selects the examples that best expose the paper's main distinctions.

\subsection{Same-Isolation Ordinary Calls versus Same-Isolation \code{sending}}

\codefile{snippets/same-vs-cross-sending.swift}

This example is the cleanest operational demonstration of the paper's thesis. The same \code{sending} parameter can be non-consuming when caller and callee share the same concrete actor isolation, yet consuming when the same value crosses to a different actor. Just as importantly, same-isolation ordinary calls and same-isolation \code{sending} calls are not identical. Ordinary calls may refine a disconnected value into an actor-bound region, whereas same-isolation \code{sending} preserves the original region. The difference is invisible at the type name level and visible only in the output environment.

\subsection{Task Capture Separates Inheritance from Transfer}

\codefile{snippets/task-capture.swift}

The closure passed to \code{Task \{ \ldots \}} is itself sent to a new task. What happens to the captured environment depends on whether the closure inherits a same concrete actor context. Actor-isolated inherited tasks preserve captures. Detached tasks and nonisolated task creation sites consume disconnected non-\code{Sendable} captures. The case study therefore repeats the same structural split already seen for values: evidence of same isolation preserves access; lack of that evidence forces transfer.

\subsection{\code{@nonisolated async} Is Caller-Inheriting, Not \code{@concurrent}}

\codefile{snippets/nonisolated-async-task-region.swift}

After SE-0461~\cite{se0461}, nonisolated async functions run on the caller's actor by default, but their non-\code{Sendable} parameters still behave like task-bound values rather than ordinary actor-bound values. This is why the parameter in the listing can be captured by a nonisolated closure and by an \code{@isolated(any)} closure, yet is rejected by a \code{@MainActor} closure. The example also illustrates a subtler point: caller inheritance does not imply that every non-\code{Sendable} result becomes disconnected. A fresh result may be transfer-ready, whereas a returned input can remain task-bound.

\subsection{Dynamic Isolation and Caller-Side Evidence}

\codefile{snippets/isolated-any-await.swift}

Dynamic isolation and explicit same-isolation evidence are mirror images of one another. A value of type \code{@isolated(any) () -> Void} carries actor identity dynamically~\cite{se0431}, so the call site must assume cross-isolation behavior and write \code{await}. By contrast, \code{isolated} parameters~\cite{se0313} and \code{\#isolation}~\cite{se0420} provide same-isolation evidence to the type checker.

\codefile{snippets/isolated-parameter.swift}

The \code{isolated} parameter example shows that no separate call rule family is needed. The passed actor witness simply determines the callee isolation, after which the ordinary same-isolation and cross-isolation rules apply. The \code{\#isolation} mechanism used in higher-order helpers follows the same principle from the caller side.

\codefile{snippets/isolation-macro.swift}

The \code{\#isolation} helper makes same-isolation reasoning first-class. Because the closure argument is checked in the caller's isolation rather than across a boundary, the closure does not need \code{@Sendable}, mutable local state remains accessible, and the captured non-\code{Sendable} value is not consumed.

\subsection{Cross-Isolation Results Are Stricter Than Parameters}

\codefile{snippets/nonsending-result.swift}

Cross-isolation parameters may transfer implicitly when the caller relinquishes ownership. Cross-isolation results are stricter because the caller is receiving a value whose region must already be known to be safe. That is why non-\code{Sendable} cross-isolation results require explicit \code{sending}, whereas disconnected arguments can be transferred implicitly. The asymmetry is one of the clearest places where the affine environment view is superior to slogan-level reasoning about isolation.

\begin{propositionbox}{Why these examples matter together}
The examples above are not independent curiosities. They all test the same joint thesis: execution capability answers where code runs, region answers where data lives, and affine update answers what survives after the interaction. Whenever a diagnosis looks surprising, one of those three questions was probably answered implicitly rather than explicitly.
\end{propositionbox}

\section{Validation Against Compiler Behavior}
\label{sec:validation}

The formalization is intentionally empirical. It takes the proposal cluster as the design background, but it treats observed compiler behavior as the calibration target. A rule that reads well yet fails against a focused probe suite is not accepted as part of the paper's account.

The validation corpus is organized around five questions.

\begin{itemize}
  \item Which annotation changes compile as direct function conversions, and which require explicit wrapping?
  \item When does same-isolation use preserve a non-\code{Sendable} value, and when does it refine or consume it?
  \item Which construction paths preserve dynamic actor identity for \code{@isolated(any)} values?
  \item When do closures keep actor-instance isolation, and when do they silently fall back to \code{@nonisolated}?
  \item How do nonisolated async calls behave on fresh versus aliased non-\code{Sendable} results after SE-0461?
\end{itemize}

\begin{definitionbox}{Kinds of evidence}
The corpus includes positive compilation tests, negative diagnostics, and property-style checks on observable values such as \code{f.isolation}. The paper relies on all three. Positive tests establish permissiveness, negative tests establish rejection boundaries, and observational tests establish which dynamic information survived a conversion path.
\end{definitionbox}

\subsection{Dynamic Isolation Identity Is Path Sensitive}

\codefile{snippets/isolated-any-tests.swift}

These tests are among the sharpest validation artifacts in the paper because they distinguish three superficially similar construction paths.

\begin{itemize}
  \item Direct conversion from \code{@MainActor () -> Void} to \code{@isolated(any) () -> Void} preserves actor identity.
  \item Building an \code{@isolated(any)} closure literal inside a \code{@MainActor} context also preserves identity through closure inference.
  \item Converting through a plain \code{() -> Void} erases actor identity, and coercing back to \code{@isolated(any)} does not recover it.
\end{itemize}

This is exactly the behavior predicted by the paper's distinction between direct coercion, contextual closure formation, and identity-erasing nonisolated conversion.

\subsection{Actor-Instance Closure Inference Is Capture Sensitive}

\codefile{snippets/actor-instance-closure.swift}

Global-actor inheritance is type-level and does not require a captured actor value. Actor-instance inheritance is different. The first closure above captures \code{self.state}, so the effective closure capability remains the actor instance. The second closure does not capture the isolated parameter \code{self}, so the compiler is free to treat it as nonisolated. This distinction validates the auxiliary predicate $\capturesIsolatedParam$ and shows why closure inference cannot be modeled as a uniform ``inherit parent isolation'' rule.

\subsection{Counterexamples Matter as Much as Success Cases}

The paper's more delicate claims are justified by negative tests and counterexamples rather than by happy-path compilation alone. Three examples are especially important.

\begin{itemize}
  \item The same-isolation and cross-isolation \code{sending} tests show that transfer is conditional on actor-equality evidence, not on the presence of the keyword alone.
  \item The nonisolated async identity counterexample shows that a returned non-\code{Sendable} value need not become disconnected merely because the callee is caller-inheriting after SE-0461. Fresh results and aliased results must be distinguished.
  \item The cross-isolation non-\code{sending} result tests show that result positions are stricter than parameter positions, validating the paper's asymmetry claim for \code{sending}.
\end{itemize}

\begin{propositionbox}{Interpretation of the evidence}
The validation results do not prove the formalization correct in the theorem-proving sense. They justify a narrower claim: within the compiler behaviors tested here, the capability-region-affine account predicts both successful and failing cases more faithfully than a one-dimensional notion of isolation.
\end{propositionbox}

This empirical stance is also a practical commitment. If future compiler behavior changes the observed conversion or transfer boundaries, the paper should change with it. The formalization is meant to remain aligned with executable evidence, not to float above it as a timeless prose summary.

\section{Discussion and Limits}
\label{sec:discussion}

The paper makes a deliberately ambitious but bounded claim. It does not claim that Swift Concurrency reduces to three words. It claims that the combination of capability, region, and affine environment update is the right explanatory vocabulary for the cluster of behaviors that programmers find hardest to predict in practice.

That claim has three strengths.

\begin{itemize}
  \item It explains why syntactically distant features, such as \code{sending}, \code{Task} capture, \code{@isolated(any)}, and \code{\#isolation}, repeatedly trigger the same kinds of surprises.
  \item It preserves the asymmetries that the compiler actually enforces, especially
    parameter-versus-result transfer and global-actor versus actor-instance closure inference.
  \item It leaves room for empirical calibration. The account is precise enough to fail against tests, which is exactly what a useful explanatory model should permit.
\end{itemize}

\begin{propositionbox}{What the model clarifies}
The model clarifies that Swift Concurrency irregularity is not primarily a keyword problem. It is a problem of conflated semantic dimensions. Surface syntax often hides whether the relevant fact is execution capability, value region, or caller-side ownership update. Once those dimensions are separated, many ``special cases'' become routine consequences of the same discipline.
\end{propositionbox}

Important limits remain.

\begin{itemize}
  \item The account is not a mechanized proof of Swift. It is a compiler-calibrated formalization.
  \item It abstracts away from many language surfaces, standard-library APIs, and optimizer-level details.
  \item It treats closure wrapping as part of the operational conversion story rather than reducing every successful conversion to one uniform semantic relation.
  \item It assumes the proposal cluster and the observed compiler behavior remain close enough that one explanatory account can serve both.
\end{itemize}

These limits are acceptable because they match the paper's purpose. The aim is not to replace the language reference or the compiler implementation. The aim is to make the existing system thinkable enough that a reader can predict new cases instead of memorizing diagnostics one by one.

\subsection{Related Work}

The capability-region decomposition in this paper is informed by several threads in the programming-language theory literature.

Region-based memory management originates with Tofte and Talpin~\cite{tofte1997region}, who showed that stack-like region lifetimes can replace garbage collection for a polymorphic lambda calculus. Grossman et~al.~\cite{grossman2002cyclone} brought regions to the systems language Cyclone, demonstrating that region disciplines can coexist with low-level programming. Swift's region-based isolation (SE-0414) inherits the core insight—tracking value connectivity by region—but applies it to concurrency rather than memory deallocation.

The static-capability approach was developed by Walker, Crary, and Morrisett~\cite{walker2000capabilities}, who used capabilities to govern access to typed memory. Chargu\'{e}raud and Pottier~\cite{chargueraud2008capabilities} further explored capability calculi with functional translations. Walker and Watkins~\cite{walker2001regions} connected regions and linear types, showing how affine or linear disciplines can track region liveness. The present paper's use of execution capability to determine \emph{where} code runs, and affine environment update to determine \emph{what} remains usable, echoes this tradition.

Most directly relevant is Milano, Turcotti, and Myers~\cite{milano2022fearless}, who proposed a flexible type system for fearless concurrency combining regions, capabilities, and linearity. Their system demonstrated that capability-region typing can ensure safe concurrent access without requiring all shared values to be deeply immutable. The lineage is not merely intellectual: Joshua Turcotti co-authored both that paper and SE-0414~\cite{se0414}, which introduced region-based isolation into Swift. Swift's approach differs in surface design—using actor isolation and \code{Sendable} rather than explicit capability annotations—but the underlying semantic decomposition is strikingly parallel.

\subsection{Future Work}

Three directions seem most promising. First, a mechanized proof in a proof assistant such as Lean or Coq would strengthen the paper's claims from compiler-calibrated observation to verified soundness. The main obstacle is faithfully modeling Swift's contextual type inference and closure formation, which interact with region tracking in ways that are easy to describe informally but nontrivial to encode. Second, the model could be extended to cover distributed actors, where network boundaries introduce latency and failure modes that the current capability-region account does not address. Third, as Swift's ownership and borrowing features mature, the interaction between move-only types, consuming parameters, and the affine environment update described here will need a unified treatment. The current model's output environment already provides a natural attachment point for ownership effects.

\section{Conclusion}
\label{sec:conclusion}

Swift Concurrency looks irregular when it is approached as a bag of diagnostics. It becomes substantially more regular when the same behaviors are grouped around three questions.

\begin{itemize}
  \item Which capability determines where this computation runs?
  \item Which region records where this non-\code{Sendable} value currently lives?
  \item How does the caller-side environment change after this interaction?
\end{itemize}

Those questions are sufficient to explain the phenomena that motivated the paper: same-isolation \code{sending} without consumption, ordinary same-isolation binding, result-side \code{sending} asymmetry, dynamic \code{@isolated(any)} behavior, task capture, actor-instance closure inference, and caller-inheriting nonisolated async execution after SE-0461~\cite{se0461}.

The main value of the paper is therefore not a new keyword or a new compiler feature. It is a sharper mental model. Capability tells the reader where code runs. Region tells the reader where data lives. Affine update tells the reader what survives. These three notions are not novel in isolation—they trace back through decades of work on region types~\cite{tofte1997region}, static capabilities~\cite{walker2000capabilities}, and their application to safe concurrency~\cite{milano2022fearless}—but their coordinated use as a lens on Swift Concurrency is new. Once those roles are separated, Swift Concurrency becomes less a collection of surprises and more a structured type discipline whose corner cases can be reasoned about in advance.


\appendix
\clearpage
\section{Notation and Auxiliary Definitions}
\label{sec:appendix-notation}

This appendix collects the auxiliary definitions that the main body uses selectively and the full rule catalog depends on throughout.

\subsection{Judgment Forms}

\[
  \Gamma;\; @\kappa;\; \alpha \vdash e : T \at \rho \dashv \Gamma'
\]

\[
  \Delta \vdash_d \mathit{decl}
\]

\begin{itemize}
  \item $\Gamma$ / $\Gamma'$: input and output environments
  \item $@\kappa$: current execution capability
  \item $\alpha$: ambient mode, either $\mathit{sync}$ or $\mathit{async}$
  \item $\rho$: result region
  \item $\Delta$: declaration environment used only to decide how bodies are checked
\end{itemize}

The declaration judgment does not produce a value type. Its role is to determine the capability and ambient mode under which a function or task body is checked.

\subsection{Capability and Function Annotations}

\[
\begin{aligned}
  @\kappa &::= @\textsf{nonisolated} \mid @\isolated(a) \\
  @\sigma &::= \cdot \mid @\textsf{Sendable} \\
  @\iota &::= @\kappa \mid @\textsf{isolated(any)} \mid @\textsf{concurrent} \\
  @\varphi &::= @\sigma\; @\iota
\end{aligned}
\]

\[
\begin{aligned}
  \toCapability(@\kappa) &= @\kappa \\
  \toCapability(@\textsf{isolated(any)}) &= @\textsf{nonisolated} \\
  \toCapability(@\textsf{concurrent}) &= @\textsf{nonisolated} \\
  \isActorIsolated(@\kappa) &\iff @\kappa \notin \{@\textsf{nonisolated}\}
\end{aligned}
\]

The split between $@\sigma$ and $@\iota$ is deliberate. \code{@Sendable} constrains capture, while $@\iota$ determines body capability. \code{@concurrent} and caller-inheriting nonisolated async functions therefore agree on body capability and diverge only at call sites.

\subsection{Regions, Normal Forms, and Merge}

\[
\begin{aligned}
  \rho_{\mathit{ns}} &::= \disconnected \mid \isolated(a) \mid \task \mid \invalid \\
  \rho &::= \rho_{\mathit{ns}} \mid \mathunderscore
\end{aligned}
\]

\[
  \forall (x : T \at \rho) \in \Gamma.\; (T : \textsf{Sendable} \Leftrightarrow \rho = \mathunderscore)
\]

$\mathunderscore$ is the normalized region tag for \code{Sendable} values. Region merge ranges only over $\rho_{\mathit{ns}}$.

\[
\begin{aligned}
  \Accessible(@\textsf{nonisolated}) &= \{\disconnected, \task, \mathunderscore\} \\
  \Accessible(@\isolated(a)) &= \{\disconnected, \isolated(a), \task, \mathunderscore\} \\
  \Capturable(@\textsf{nonisolated}) &= \{\disconnected, \task, \mathunderscore\} \\
  \Capturable(@\isolated(a)) &= \{\disconnected, \isolated(a), \mathunderscore\}
\end{aligned}
\]

The difference between $\Accessible$ and $\Capturable$ is concentrated in $\task$: task-bound values may be directly readable in some actor-isolated contexts but may not be captured into actor-isolated closures.

\[
\begin{aligned}
  \toRegion(@\isolated(a)) &= \isolated(a) \\
  \toRegion(@\textsf{nonisolated}) &= \disconnected
\end{aligned}
\]

\[
\begin{aligned}
  \disconnected \sqcup \disconnected &= \disconnected \\
  \disconnected \sqcup \isolated(a) &= \isolated(a) \\
  \disconnected \sqcup \task &= \task \\
  \isolated(a) \sqcup \isolated(a) &= \isolated(a) \\
  \isolated(a) \sqcup \isolated(b) &= \invalid \text{ for } a \neq b \\
  \isolated(a) \sqcup \task &= \invalid \\
  \task \sqcup \task &= \task
\end{aligned}
\]

\subsection{Closure Support}

\[
  \isAllSendable(\Gamma_{\mathit{captured}})
    \iff \forall (x : T \at \rho) \in \Gamma_{\mathit{captured}}.\; T : \textsf{Sendable}
\]

\[
\begin{aligned}
  \inheritActorContext(\mathit{closure}) &: \mathit{Bool} \\
  \isPassedToSendingParameter(\mathit{closure}) &: \mathit{Bool} \\
  \capturesIsolatedParam(\mathit{closure}) &: \mathit{Bool}
\end{aligned}
\]

\[
\begin{aligned}
  \effectiveClosureCapability(@\textsf{nonisolated}, \mathit{closure}) &= @\textsf{nonisolated} \\
  \effectiveClosureCapability(@\isolated(\mathit{globalActor}), \mathit{closure}) &= @\isolated(\mathit{globalActor})
\end{aligned}
\]

\[
\begin{aligned}
  \mathit{ECC}_{\mathit{local}}(\mathit{closure})
    &= @\isolated(\mathit{localActor})
      \quad \text{if } \capturesIsolatedParam(\mathit{closure}) \\
  \mathit{ECC}_{\mathit{local}}(\mathit{closure})
    &= @\textsf{nonisolated}
      \quad \text{otherwise}
\end{aligned}
\]

Here, $\mathit{ECC}_{\mathit{local}}(\mathit{closure})$ abbreviates $\effectiveClosureCapability(@\isolated(\mathit{localActor}), \mathit{closure})$. The final clause is the critical actor-instance branch: local actor isolation is preserved only when the isolated parameter or \code{self} is actually captured.

\subsection{Transfer and Sendable Inference Support}

\[
\begin{aligned}
  [\sending] &::= \sending \mid \cdot \\
  \canSend(@\kappa, @\iota, [\sending]) &\iff [\sending] \in \{\sending\} \lor (@\kappa \neq @\iota \land @\iota \notin \{@\textsf{nonisolated}\})
\end{aligned}
\]

This predicate unifies explicit \code{sending} and implicit cross-isolation transfer.

\[
\begin{aligned}
  \isNonLocal(f) &: \mathit{Bool} \\
  \instanceMethods(T) &: \mathit{Set}\langle\mathit{MethodDecl}\rangle \\
  \hasIsolatedKeyPathComponent(\mathit{kp}) &: \mathit{Bool} \\
  \isAllSendable(\mathit{captures}(\mathit{kp})) &: \mathit{Bool}
\end{aligned}
\]

These definitions support the Sendable-inference rules for top-level functions, static methods, unapplied instance-method references, and key-path literals.

\clearpage
\section{Boundary and Relation Rules}
\label{sec:appendix-relation}

This appendix prints the non-expression rules that determine body capability, ambient async mode, and relation-level conversion structure.

\subsection{Sync and Async Boundaries}

\subsubsection{\textsc{decl-fun}}

\typerule{decl-fun}
  {\Gamma_{\mathit{body}} = x_1 : A_1 \at \rho_1, \ldots, x_n : A_n \at \rho_n \\
   @\kappa = \toCapability(\mathrm{proj}_\iota(@\varphi)) \\
   \Gamma_{\mathit{body}};\; @\kappa;\; \alpha \vdash \mathit{body} : R \at \rho \dashv \Gamma'_{\mathit{body}}}
  {\Delta \vdash_d @\varphi\; \mathbf{func}\; \mathit{foo}(x_1 : A_1, \ldots, x_n : A_n)\; \alpha \to R\; \{\mathit{body}\}}

For nonisolated bodies (sync and async alike), non-\code{Sendable} parameters are initialized in region $\task$. The $\task$ region represents caller-owned data not bound to any specific actor; even in the sync case, the function may be called from any isolation context, so assigning parameters to a specific actor domain is not statically safe. This is the source of the later capture restrictions discussed in the main body.

\subsubsection{\textsc{decl-fun-isolated-param}}

\typerule{decl-fun-isolated-param}
  {a : \mathit{ActorType},\; \Gamma_{\mathit{body}} = a : \mathit{ActorType} \at \mathunderscore,\; x_1 : A_1 \at \rho_1, \ldots, x_n : A_n \at \rho_n \\
   \Gamma_{\mathit{body}};\; @\isolated(a);\; \alpha \vdash \mathit{body} : R \at \rho \dashv \Gamma'_{\mathit{body}}}
  {\Delta \vdash_d \mathbf{func}\; \mathit{foo}(\mathit{actor} : \mathbf{isolated}\; \mathit{ActorType}, x_1 : A_1, \ldots)\; \alpha \to R\; \{\mathit{body}\}}

\subsubsection{\textsc{decl-fun-isolation-inheriting}}

\typerule{decl-fun-isolation-inheriting}
  {\Gamma_{\mathit{body}} = \mathit{iso} : (\mathbf{any}\; \mathit{Actor})? \at \mathunderscore,\; x_1 : A_1 \at \rho_1, \ldots, x_n : A_n \at \rho_n \\
   \Gamma_{\mathit{body}};\; @\textsf{nonisolated};\; \alpha \vdash \mathit{body} : R \at \rho \dashv \Gamma'_{\mathit{body}}}
  {\Delta \vdash_d \mathbf{func}\; \mathit{foo}(\mathit{isolation} : \mathbf{isolated}\; (\mathbf{any}\; \mathit{Actor})? = \texttt{\#isolation}, x_1 : A_1, \ldots)\; \alpha \to R\; \{\mathit{body}\}}

\subsubsection{\code{Task} and \code{Task.detached} bodies}

\code{Task} bodies are checked under $\alpha = \mathit{async}$. Their closure parameters are also \code{sending}, so closure capture follows the same \textsc{closure-sending-*} rules printed in Appendix~\ref{sec:appendix-full}. \code{Task.init} differs from \code{Task.detached} because the former may inherit caller actor context through \code{@\_inheritActorContext}, while the latter never does.

\subsection{Relation Rules}

The formalization distinguishes semantic subtyping from one-step contextual coercion.

\begin{itemize}
  \item $\isoSubtyping$: a semantic relation that may compose transitively
  \item $\isoCoercion$: a one-step direct coercion judgment used by \textsc{func-conv}
\end{itemize}

\subsubsection{Isolation Subtyping}

\paragraph{\textsc{iso-subtyping-identity}}
\[
  \isoSubtyping(@\sigma, \iota, @\sigma, \iota, \alpha')
\]

\paragraph{\textsc{iso-subtyping-sendable-forget}}
\[
  \isoSubtyping(@\textsf{Sendable}, \iota, \cdot, \iota, \alpha')
\]

\paragraph{\textsc{iso-subtyping-nonisolated-to-isolated-any}}
\[
  \isoSubtyping(@\sigma, @\textsf{nonisolated}, @\sigma, @\textsf{isolated(any)}, \alpha')
\]

\paragraph{\textsc{iso-subtyping-mainactor-to-isolated-any}}
\[
  \isoSubtyping(@\sigma, @\textsf{MainActor}, @\sigma, @\textsf{isolated(any)}, \alpha')
\]

\paragraph{\textsc{iso-subtyping-mainactor-implicit-sendable}}
\[
  \isoSubtyping(\cdot, @\textsf{MainActor}, @\textsf{Sendable}, @\textsf{MainActor}, \alpha')
\]

\paragraph{\textsc{iso-subtyping-isolated-local-actor-identity}}
\[
  \isoSubtyping(@\sigma, @\isolated(\mathit{localActor}), @\sigma, @\isolated(\mathit{localActor}), \alpha')
\]

\paragraph{\textsc{iso-subtyping-isolated-local-actor-sendable-forget}}
\[
  \isoSubtyping(@\textsf{Sendable}, @\isolated(\mathit{localActor}), \cdot, @\isolated(\mathit{localActor}), \alpha')
\]

\paragraph{\textsc{iso-subtyping-transitive}}

\typerule{iso-subtyping-transitive}
  {\isoSubtyping(@\sigma_1, \iota_1, @\sigma_2, \iota_2, \alpha') \\
   \isoSubtyping(@\sigma_2, \iota_2, @\sigma_3, \iota_3, \alpha')}
  {\isoSubtyping(@\sigma_1, \iota_1, @\sigma_3, \iota_3, \alpha')}

\paragraph{\textsc{iso-subtyping-to-coercion}}

\typerule{iso-subtyping-to-coercion}
  {\isoSubtyping(@\sigma_1, \iota_1, @\sigma_2, \iota_2, \alpha')}
  {\isoCoercion(@\sigma_1, \iota_1, @\sigma_2, \iota_2, \alpha')}

\subsubsection{Isolation Coercion}

\paragraph{\textsc{iso-coercion-isolated-any-to-nonisolated-deprecated}}
\[
  \isoCoercion(@\sigma, @\textsf{isolated(any)}, @\sigma, @\textsf{nonisolated}, \mathit{sync})
\]

\paragraph{\textsc{iso-coercion-sendable-nonisolated-universal}}

\typerule{iso-coercion-sendable-nonisolated-universal}
  {\iota_2 \neq @\isolated(\mathit{localActor})}
  {\isoCoercion(@\textsf{Sendable}, @\textsf{nonisolated}, @\sigma_2, \iota_2, \alpha')}

\paragraph{\textsc{iso-coercion-async-mainactor-universal}}

\typerule{iso-coercion-async-mainactor-universal}
  {\iota_2 \neq @\isolated(\mathit{localActor})}
  {\isoCoercion(@\sigma, @\textsf{MainActor}, @\sigma_2, \iota_2, \mathit{async})}

\paragraph{\textsc{iso-coercion-async-nonsendable-equiv}}

\typerule{iso-coercion-async-nonsendable-equiv}
  {\iota_1 \in \{@\textsf{nonisolated}, @\textsf{concurrent}, @\textsf{isolated(any)}\} \\
   \iota_2 \in \{@\textsf{nonisolated}, @\textsf{concurrent}, @\textsf{isolated(any)}\}}
  {\isoCoercion(\cdot, \iota_1, \cdot, \iota_2, \mathit{async})}

The conversion matrices in the source formalization remain relevant because not every successful local conversion is a one-step coercion judgment. Actor-parameter branches and several \code{@isolated(any)}-to-actor conversions require explicit closure wrapping and therefore belong to the term-level conversion story rather than to direct relation rules alone.

\clearpage
\section{Full Typing-Rule Catalog}
\label{sec:appendix-full}

This appendix prints the complete typing-rule inventory used by the paper. The main body emphasizes only a kernel, but no rule family is omitted here.

\subsection{Variables}

\subsubsection{\textsc{var}}

\typerule{var}
  {x : T \at \rho \in \Gamma \\ \rho \in \Accessible(@\kappa)}
  {\Gamma;\; @\kappa;\; \alpha \vdash x : T \at \rho \dashv \Gamma}

\subsection{Sequencing}

\subsubsection{\textsc{seq}}

\typerule{seq}
  {\Gamma;\; @\kappa;\; \alpha \vdash e_1 : () \at \rho_1 \dashv \Gamma_1 \\
   \Gamma_1;\; @\kappa;\; \alpha \vdash e_2 : T \at \rho_2 \dashv \Gamma_2}
  {\Gamma;\; @\kappa;\; \alpha \vdash e_1;\; e_2 : T \at \rho_2 \dashv \Gamma_2}

\subsection{Function Type Conversion}

\subsubsection{\textsc{sync-to-async}}

\typerule{sync-to-async}
  {\Gamma;\; @\kappa;\; \alpha \vdash f : @\varphi\; () \to B \at \rho_f \dashv \Gamma'}
  {\Gamma;\; @\kappa;\; \alpha \vdash f : @\varphi\; ()\; \mathit{async} \to B \at \rho_f \dashv \Gamma'}

\subsubsection{\textsc{func-conv}}

\typerule{func-conv}
  {\Gamma;\; @\kappa;\; \alpha \vdash f : @\sigma_1\; @\iota_1\; (A)\; \alpha' \to B \at \rho_f \dashv \Gamma' \\
   \isoCoercion(@\sigma_1, @\iota_1, @\sigma_2, @\iota_2, \alpha') \\
   \rho' = (\text{if } @\sigma_2 = @\textsf{Sendable} \text{ then } \mathunderscore \text{ else } \rho_f)}
  {\Gamma;\; @\kappa;\; \alpha \vdash f : @\sigma_2\; @\iota_2\; (A)\; \alpha' \to B \at \rho' \dashv \Gamma'}

\subsubsection{\textsc{isolated-any-to-async}}

\typerule{isolated-any-to-async}
  {\Gamma;\; @\kappa;\; \alpha \vdash f : @\sigma\; @\textsf{isolated(any)}\; () \to B \at \rho_f \dashv \Gamma' \\
   \neg\isActorIsolated(@\iota_2) \\
   \rho' = (\text{if } @\sigma = @\textsf{Sendable} \text{ then } \mathunderscore \text{ else } \rho_f)}
  {\Gamma;\; @\kappa;\; \alpha \vdash f : @\sigma\; @\iota_2\; ()\; \mathit{async} \to B \at \rho' \dashv \Gamma'}

\subsubsection{Subtyping / Coercion Validation Table}

The source account includes combined sync and async validation tables that explain which annotation changes are justified by direct coercion and which require closure wrapping. In this paper, those tables are conceptually split between Appendix~\ref{sec:appendix-relation} and the note below: Appendix~\ref{sec:appendix-relation} records the relation-level direct edges, while the present appendix records the term-level fact that some successful local conversions still require a freshly constructed closure.

\subsubsection{Conversion by Closure Wrapping}

Not every successful local conversion is a one-step coercion judgment. Some require an explicit new closure, for example when an actor parameter must be introduced or when a dynamic-isolation value must be called under \code{await} inside a freshly constructed closure. This note is part of the full function-conversion story and is therefore recorded explicitly here.

\subsection{Region Merge}

\subsubsection{\textsc{region-merge}}

\typerule{region-merge}
  {\Gamma;\; @\kappa;\; \alpha \vdash e_1 : T_1 \at \rho_1 \dashv \Gamma_1 \\
   \Gamma_1;\; @\kappa;\; \alpha \vdash e_2 : T_2 \at \rho_2 \dashv \Gamma_2 \\
   T_1 : {\sim}\textsf{Sendable},\; T_2 : {\sim}\textsf{Sendable} \\
   \rho_1, \rho_2 \in \rho_{\mathit{ns}} \\
   \rho = \rho_1 \sqcup \rho_2}
  {\Gamma;\; @\kappa;\; \alpha \vdash (e_1.\mathit{field} = e_2) : () \at \mathunderscore \dashv \Gamma_2[\rho_1 \mapsto \rho, \rho_2 \mapsto \rho]}

\subsection{Calls}

The call rules partition into four families: same-isolation, cross-isolation, explicit \code{@concurrent}, and caller-inheriting \code{@nonisolated async}. The rules below are printed in that order.

\subsubsection{\textsc{call-nonsendable-noconsume}}

\typerule{call-nonsendable-noconsume}
  {\Gamma;\; @\kappa;\; \alpha \vdash f : @\sigma\; @\iota\; ([\sending]\; A)\; \alpha' \to B \at \rho_f \dashv \Gamma_1 \\
   \Gamma_1;\; @\kappa;\; \alpha \vdash \mathit{arg} : A \at \rho_a \dashv \Gamma_2 \\
   A : {\sim}\textsf{Sendable},\; B : \textsf{Sendable},\; @\iota \notin \{@\textsf{concurrent}\} \\
   \canSend(@\kappa, @\iota, [\sending]) \\
   \rho_a \in \Accessible(@\kappa) \\
   \isActorIsolated(@\kappa) \land @\kappa = @\iota}
  {\Gamma;\; @\kappa;\; \alpha \vdash [\code{await}]\; f(\mathit{arg}) : B \at \mathunderscore \dashv \Gamma_2}

\subsubsection{\textsc{call-nonsendable-consume}}

\typerule{call-nonsendable-consume}
  {\Gamma;\; @\kappa;\; \alpha \vdash f : @\sigma\; @\iota\; ([\sending]\; A)\; \alpha' \to B \at \rho_f \dashv \Gamma_1 \\
   \Gamma_1;\; @\kappa;\; \alpha \vdash \mathit{arg} : A \at \disconnected \dashv \Gamma_2 \\
   A : {\sim}\textsf{Sendable},\; B : \textsf{Sendable},\; @\iota \notin \{@\textsf{concurrent}\} \\
   \canSend(@\kappa, @\iota, [\sending]) \\
   \neg(\isActorIsolated(@\kappa) \land @\kappa = @\iota)}
  {\Gamma;\; @\kappa;\; \alpha \vdash [\code{await}]\; f(\mathit{arg}) : B \at \mathunderscore \dashv (\Gamma_2 \setminus \{\mathit{arg}\})}

\subsubsection{\textsc{call-same-sync-sendable}}

\typerule{call-same-sync-sendable}
  {\Gamma;\; @\kappa;\; \alpha \vdash f : @\sigma\; @\kappa\; (A) \to B \at \rho_f \dashv \Gamma_1 \\
   \Gamma_1;\; @\kappa;\; \alpha \vdash \mathit{arg} : A \at \mathunderscore \dashv \Gamma_2 \\
   \rho \in \Accessible(@\kappa)}
  {\Gamma;\; @\kappa;\; \alpha \vdash f(\mathit{arg}) : B \at \rho \dashv \Gamma_2}

\subsubsection{\textsc{call-same-async-sendable}}

\typerule{call-same-async-sendable}
  {\Gamma;\; @\kappa;\; \mathit{async} \vdash f : @\sigma\; @\kappa\; (A)\; \mathit{async} \to B \at \rho_f \dashv \Gamma_1 \\
   \Gamma_1;\; @\kappa;\; \mathit{async} \vdash \mathit{arg} : A \at \mathunderscore \dashv \Gamma_2 \\
   A : \textsf{Sendable} \\
   \rho \in \Accessible(@\kappa)}
  {\Gamma;\; @\kappa;\; \mathit{async} \vdash \mathbf{await}\; f(\mathit{arg}) : B \at \rho \dashv \Gamma_2}

\subsubsection{\textsc{call-same-nonsendable-merge}}

\typerule{call-same-nonsendable-merge}
  {\Gamma;\; @\kappa;\; \alpha \vdash f : @\sigma\; @\kappa\; (A) \to B \at \rho_f \dashv \Gamma_1 \\
   \Gamma_1;\; @\kappa;\; \alpha \vdash \mathit{arg} : A \at \rho_a \dashv \Gamma_2 \\
   A : {\sim}\textsf{Sendable} \\
   \rho_a \in \Accessible(@\kappa) \\
   \rho_{a'} = \rho_a \sqcup \toRegion(@\kappa) \\
   \rho_b \in \Accessible(@\kappa),\; (B : \textsf{Sendable} \Rightarrow \rho_b = \mathunderscore)}
  {\Gamma;\; @\kappa;\; \alpha \vdash f(\mathit{arg}) : B \at \rho_b \dashv (\Gamma_2[\mathit{arg} \mapsto A \at \rho_{a'}])}

\subsubsection{\textsc{call-cross-sendable}}

\typerule{call-cross-sendable}
  {\Gamma;\; @\kappa;\; \mathit{async} \vdash f : @\sigma\; @\iota\; (A)\; \alpha \to B \at \rho_f \dashv \Gamma_1 \\
   \Gamma_1;\; @\kappa;\; \mathit{async} \vdash \mathit{arg} : A \at \mathunderscore \dashv \Gamma_2 \\
   A : \textsf{Sendable},\; B : \textsf{Sendable},\; @\kappa \neq @\iota,\; @\iota \notin \{@\textsf{concurrent}, @\textsf{nonisolated}\}}
  {\Gamma;\; @\kappa;\; \mathit{async} \vdash \mathbf{await}\; f(\mathit{arg}) : B \at \mathunderscore \dashv \Gamma_2}

\subsubsection{\textsc{call-cross-sending-result}}

\typerule{call-cross-sending-result}
  {\Gamma;\; @\kappa;\; \mathit{async} \vdash f : @\sigma\; @\iota\; ()\; \alpha \to \sending\; B \at \rho_f \dashv \Gamma' \\
   B : {\sim}\textsf{Sendable},\; @\kappa \neq @\iota,\; @\iota \notin \{@\textsf{concurrent}\}}
  {\Gamma;\; @\kappa;\; \mathit{async} \vdash \mathbf{await}\; f() : B \at \disconnected \dashv \Gamma'}

\subsubsection{\textsc{call-cross-nonsending-result-error}}

\typerule{call-cross-nonsending-result-error}
  {\Gamma;\; @\kappa;\; \mathit{async} \vdash f : @\sigma\; @\iota\; ()\; \alpha \to B \at \rho_f \dashv \Gamma_1 \\
   B : {\sim}\textsf{Sendable},\; @\kappa \neq @\iota,\; @\iota \notin \{@\textsf{concurrent}\}}
  {\text{derivation fails (compile error)}}

\subsubsection{\textsc{call-concurrent-sendable}}

\typerule{call-concurrent-sendable}
  {\Gamma;\; @\kappa;\; \mathit{async} \vdash f : @\sigma\; @\textsf{concurrent}\; (A)\; \mathit{async} \to B \at \rho_f \dashv \Gamma_1 \\
   \Gamma_1;\; @\kappa;\; \mathit{async} \vdash \mathit{arg} : A \at \mathunderscore \dashv \Gamma_2 \\
   A : \textsf{Sendable},\; B : \textsf{Sendable}}
  {\Gamma;\; @\kappa;\; \mathit{async} \vdash \mathbf{await}\; f(\mathit{arg}) : B \at \mathunderscore \dashv \Gamma_2}

\subsubsection{\textsc{call-concurrent-nonsendable}}

\typerule{call-concurrent-nonsendable}
  {\Gamma;\; @\kappa;\; \mathit{async} \vdash f : @\sigma\; @\textsf{concurrent}\; (A)\; \mathit{async} \to B \at \rho_f \dashv \Gamma_1 \\
   \Gamma_1;\; @\kappa;\; \mathit{async} \vdash \mathit{arg} : A \at \disconnected \dashv \Gamma_2 \\
   A : {\sim}\textsf{Sendable},\; B : \textsf{Sendable}}
  {\Gamma;\; @\kappa;\; \mathit{async} \vdash \mathbf{await}\; f(\mathit{arg}) : B \at \mathunderscore \dashv \Gamma_2}

\subsubsection{\textsc{call-nonisolated-async-inherit}}

\typerule{call-nonisolated-async-inherit}
  {\Gamma;\; @\kappa;\; \mathit{async} \vdash f : @\textsf{nonisolated}\; (A)\; \mathit{async} \to B \at \rho_f \dashv \Gamma_1 \\
   \Gamma_1;\; @\kappa;\; \mathit{async} \vdash \mathit{arg} : A \at \rho_a \dashv \Gamma_2}
  {\Gamma;\; @\kappa;\; \mathit{async} \vdash \mathbf{await}\; f(\mathit{arg}) : B \at \rho_b \dashv \Gamma_2}

The unconstrained $\rho_b$ is intentional. A caller-inheriting nonisolated async function may return a fresh disconnected value or may return an alias of a task-bound input.

\subsubsection{\textsc{call-nonisolated-sync}}

\typerule{call-nonisolated-sync}
  {\Gamma;\; @\kappa;\; \alpha \vdash f : @\sigma\; @\textsf{nonisolated}\; (A) \to B \at \rho_f \dashv \Gamma_1 \\
   \Gamma_1;\; @\kappa;\; \alpha \vdash \mathit{arg} : A \at \rho_a \dashv \Gamma_2 \\
   @\kappa \neq @\textsf{nonisolated} \\
   \rho_a \in \Accessible(@\kappa) \\
   \rho_b \in \Accessible(@\kappa),\; (B : \textsf{Sendable} \Rightarrow \rho_b = \mathunderscore)}
  {\Gamma;\; @\kappa;\; \alpha \vdash f(\mathit{arg}) : B \at \rho_b \dashv \Gamma_2}

A nonisolated synchronous function called from an actor-isolated context executes within the caller's isolation domain without \code{await}. Neither argument consumption nor environment binding occurs. The result region~$\rho_b$ follows the same generalization as \textsc{call-same-nonsendable-merge}: \code{Sendable} results normalize to~$\mathunderscore$; non-\code{Sendable} results may be disconnected or actor-bound.

\subsubsection{Explanatory Notes}

\paragraph{\textsc{call-isolated-param-semantics}}
Functions with an \code{isolated} parameter do not require a separate family of call rules. Instead, the passed actor instance determines the callee isolation, after which the existing same-isolation and cross-isolation rules apply in the ordinary way.

\paragraph{\textsc{call-isolation-macro-semantics}}
Calls that use \code{isolated (any Actor)? = \#isolation} are interpreted with a caller-provided same-isolation witness. This is why \code{\#isolation} enables non-\code{@Sendable} closure capture and mutable local access without introducing a cross-isolation boundary.

\subsection{\code{@isolated(any)}}

\subsubsection{\textsc{isolated-any-isolation-prop} (\code{f.isolation})}

\typerule{isolated-any-isolation-prop}
  {\Gamma;\; @\kappa;\; \alpha \vdash f : @\textsf{isolated(any)}\; () \to B \at \rho_f \dashv \Gamma'}
  {\Gamma;\; @\kappa;\; \alpha \vdash f.\mathit{isolation} : (\mathbf{any}\; \mathit{Actor})? \at \mathunderscore \dashv \Gamma'}

This rule exposes the dynamic actor identity carried by an \code{@isolated(any)} value. Direct actor-aware construction paths may preserve that identity, while conversion through an ordinary nonisolated function type may erase it permanently.

\subsection{Closures}

\subsubsection{\textsc{closure-inherit-parent}}

\typerule{closure-inherit-parent}
  {@\kappa_{\mathit{eff}} = \effectiveClosureCapability(@\kappa, \{e\}) \\
   \Gamma_{\mathit{captured}} \subseteq \Gamma \land \neg\isPassedToSendingParameter(\{e\}) \\
   \forall (y : U \at \rho) \in \Gamma_{\mathit{captured}},\; \rho \in \Capturable(@\kappa_{\mathit{eff}}) \\
   \Gamma_{\mathit{captured}};\; @\kappa_{\mathit{eff}};\; \alpha' \vdash e : B \at \rho_{\mathit{ret}} \dashv \Gamma'_{\mathit{cl}} \\
   \rho_{\mathit{closure}} = \bigsqcup\{\rho \mid (y : T \at \rho) \in \Gamma_{\mathit{captured}} \land T : {\sim}\textsf{Sendable}\}}
  {\Gamma;\; @\kappa;\; \alpha \vdash \{e\} : @\kappa_{\mathit{eff}}\; ()\; \alpha' \to B \at \rho_{\mathit{closure}} \dashv \Gamma}

For global actors, $\effectiveClosureCapability$ keeps the inherited actor capability without requiring a captured actor value. For actor instances, it preserves instance isolation only when the isolated parameter or \code{self} is actually captured.

\subsubsection{\textsc{closure-no-inherit-parent}}

\typerule{closure-no-inherit-parent}
  {\Gamma_{\mathit{captured}} \subseteq \Gamma \\
   \neg\inheritActorContext(\{e\}) \land \neg\isPassedToSendingParameter(\{e\}) \\
   \isAllSendable(\Gamma_{\mathit{captured}}) \\
   @\kappa_{\mathit{cl}} = \toCapability(@\iota) \\
   \Gamma_{\mathit{captured}};\; @\kappa_{\mathit{cl}};\; \alpha' \vdash e : B \at \rho_{\mathit{ret}} \dashv \Gamma'_{\mathit{cl}}}
  {\Gamma;\; @\kappa;\; \alpha \vdash \{e\} : @\textsf{Sendable}\; @\iota\; ()\; \alpha' \to B \at \mathunderscore \dashv \Gamma}

\subsubsection{\textsc{closure-sending-noconsume}}

\typerule{closure-sending-noconsume}
  {f : (\sending(@\kappa_{\mathit{cl}}\; ()\; \alpha' \to T)) \to R \\
   \Gamma_{\mathit{captured}} \subseteq \Gamma \\
   \inheritActorContext(\{e\}) \land \isActorIsolated(@\kappa) \land @\kappa_{\mathit{cl}} = @\kappa \\
   \forall (x : U \at \rho) \in \Gamma_{\mathit{captured}},\; \rho \in \Capturable(@\kappa) \\
   \Gamma_{\mathit{captured}};\; @\kappa_{\mathit{cl}};\; \alpha' \vdash e : T \at \rho_{\mathit{ret}} \dashv \Gamma'_{\mathit{cl}}}
  {\Gamma;\; @\kappa;\; \alpha \vdash f(\{e\}) : R \at \mathunderscore \dashv \Gamma}

\subsubsection{\textsc{closure-sending-consume}}

\typerule{closure-sending-consume}
  {f : (\sending(@\kappa_{\mathit{cl}}\; ()\; \alpha' \to T)) \to R \\
   \Gamma_{\mathit{captured}} \subseteq \Gamma \\
   \neg(\inheritActorContext(\{e\}) \land \isActorIsolated(@\kappa) \land @\kappa_{\mathit{cl}} = @\kappa) \\
   \forall (x : U \at \rho) \in \Gamma_{\mathit{captured}}.\; (U : \textsf{Sendable}) \lor (U : {\sim}\textsf{Sendable} \land \rho = \disconnected) \\
   \Gamma_{\mathit{captured}};\; @\kappa_{\mathit{cl}};\; \alpha' \vdash e : T \at \rho_{\mathit{ret}} \dashv \Gamma'_{\mathit{cl}} \\
   \Gamma' = \Gamma \setminus \{x \mid (x : U \at \disconnected) \in \Gamma_{\mathit{captured}} \land U : {\sim}\textsf{Sendable}\}}
  {\Gamma;\; @\kappa;\; \alpha \vdash f(\{e\}) : R \at \mathunderscore \dashv \Gamma'}

This rule is the transfer-side counterpart of \textsc{call-nonsendable-consume}: the sent closure leaves the caller, so disconnected non-\code{Sendable} captures are removed from the caller environment unless same concrete actor inheritance proves the non-consuming branch.

\subsection{Async Let}

\subsubsection{\textsc{async-let}}

\typerule{async-let}
  {\Gamma_{\mathit{captured}} \subseteq \Gamma \\
   \forall (y : U \at \rho) \in \Gamma_{\mathit{captured}}.\; (U : \textsf{Sendable}) \lor (U : {\sim}\textsf{Sendable} \land \rho = \disconnected) \\
   \Gamma_{\mathit{captured}};\; @\textsf{nonisolated};\; \mathit{async} \vdash \mathit{expr} : T \at \rho_{\mathit{init}} \dashv \Gamma'_{\mathit{child}} \\
   \Gamma_{\mathit{sent}} = \Gamma \setminus \{y \mid (y : U \at \disconnected) \in \Gamma_{\mathit{captured}} \land U : {\sim}\textsf{Sendable}\} \\
   \rho_x = (\text{if } T : \textsf{Sendable} \text{ then } \mathunderscore \text{ else } \disconnected)}
  {\Gamma;\; @\kappa;\; \mathit{async} \vdash \mathbf{async\; let}\; x = \mathit{expr} \dashv \Gamma_{\mathit{sent}},\; x : T \at \rho_x}

The initializer is wrapped in an implicit nonisolated async autoclosure that forms a child task boundary. Because the conclusion already requires ambient mode~$\mathit{async}$, \code{async let} is available only in async contexts. Non-\code{Sendable} captures must be disconnected and are temporarily consumed. The output environment~$\Gamma'_{\mathit{child}}$ records which captures survive cross-isolation calls within the child task body.

\subsubsection{\textsc{async-let-access}}

\typerule{async-let-access}
  {(x : T \at \rho) \in \Gamma \quad \isAsyncLet(x) \\
   \Gamma' = \undosend(\Gamma, x)}
  {\Gamma;\; @\kappa;\; \mathit{async} \vdash \mathbf{await}\; x : T \at \rho \dashv \Gamma'}

Accessing an \code{async let} binding requires \code{await} (and \code{try} if the initializer can throw). The $\undosend$ function restores captures that were not consumed by cross-isolation within the child task: $\undosend(\Gamma, x) = \Gamma \cup \{y : U \at \disconnected \mid y \in \mathit{sent}_{\text{async-let}}(x) \land y \in \mathrm{dom}(\Gamma'_{\mathit{child}})\}$. Here $\mathrm{dom}(\Gamma)$ denotes the domain of environment~$\Gamma$, i.e.\ the set of variable names $y$ such that $\exists T, \rho.\; y : T \at \rho \in \Gamma$. This is what distinguishes \code{async let} from permanent \code{sending}-based transfer.

\subsection{Sendable Inference}

\subsubsection{\textsc{infer-sendable-nonlocal}}

\typerule{infer-sendable-nonlocal}
  {\isNonLocal(f) \\ f : (A_1, \ldots, A_n)\; \alpha \to R}
  {\Gamma;\; @\kappa;\; \alpha' \vdash f : @\textsf{Sendable}\; (A_1, \ldots, A_n)\; \alpha \to R \at \mathunderscore \dashv \Gamma}

\subsubsection{\textsc{infer-sendable-method-sendable}}

\typerule{infer-sendable-method-sendable}
  {T : \textsf{Sendable} \\ m \in \instanceMethods(T) \\ m : (A_1, \ldots, A_n)\; \alpha \to R}
  {\Gamma;\; @\kappa;\; \alpha' \vdash T.m : @\textsf{Sendable}\; (T) \to @\textsf{Sendable}\; ((A_1, \ldots, A_n)\; \alpha \to R) \at \mathunderscore \dashv \Gamma}

\subsubsection{\textsc{infer-sendable-method-non-sendable}}

\typerule{infer-sendable-method-non-sendable}
  {T : {\sim}\textsf{Sendable} \\ m \in \instanceMethods(T) \\ m : (A_1, \ldots, A_n)\; \alpha \to R}
  {\Gamma;\; @\kappa;\; \alpha' \vdash T.m : @\textsf{Sendable}\; (T) \to ((A_1, \ldots, A_n)\; \alpha \to R) \at \mathunderscore \dashv \Gamma}

\subsubsection{\textsc{infer-sendable-keypath}}

\typerule{infer-sendable-keypath}
  {\mathit{kp} = \mathit{KeyPath}(T, \mathit{path}) \\
   \neg\hasIsolatedKeyPathComponent(\mathit{kp}) \\
   \isAllSendable(\mathit{captures}(\mathit{kp}))}
  {\Gamma;\; @\kappa;\; \alpha \vdash \mathit{kp} : \mathbf{any}\; \mathit{KeyPath}\langle T, V\rangle\; \&\; \textsf{Sendable} \at \mathunderscore \dashv \Gamma}

\clearpage
\section{Complete Example Inventory}
\label{sec:appendix-examples}

This appendix preserves the full worked-example set in compact paper form. The main body discusses only the most revealing cases in depth.

\subsection{Example 1. Same-isolation ordinary use binds instead of consuming}

\textsc{call-same-nonsendable-merge} refines the caller environment. A disconnected value passed to an ordinary same-isolation non-\code{sending} parameter remains usable, but it may become actor-bound and therefore lose eligibility for later transfer.

\subsection{Example 2. Same-isolation \code{sending} versus cross-isolation consume}

When $\isActorIsolated(@\kappa)$ and $@\kappa = @\iota$ hold, \textsc{call-nonsendable-noconsume} preserves the original region. When that same-isolation evidence fails, \textsc{call-nonsendable-consume} removes the binding from the caller environment.

\codefile{snippets/same-vs-cross-sending.swift}

\subsection{Example 3. Cross-isolation implicit transfer consumes}

Explicit \code{sending} is not required for transfer. If a non-\code{Sendable} argument crosses to a distinct actor-isolated callee, $\canSend$ holds implicitly and the caller binding is consumed.

\subsection{Example 4. Cross-isolation explicit \code{sending} consumes}

An explicit \code{sending} parameter across isolation boundaries is the direct transfer case: the caller must provide a disconnected value and loses the binding after the call.

\subsection{Example 5. \code{Task} capture distinguishes inherited and transferred closures}

\code{Task.init} may preserve captures when same concrete actor inheritance is available.
\code{Task.detached} and nonisolated task creation fall into the transfer branch and consume disconnected non-\code{Sendable} captures.

\codefile{snippets/task-capture.swift}

\subsection{Example 6. Nonisolated parameters behave like task regions}

Non-\code{Sendable} parameters of a nonisolated function (sync or async) are task-bound rather than ordinary actor-bound values. This explains the closure-capture asymmetry between nonisolated or \code{@isolated(any)} closures and \code{@MainActor} closures after SE-0461.

\codefile{snippets/nonisolated-async-task-region.swift}

\subsection{Example 7. \code{@isolated(any)} requires \code{await}}

Dynamic actor identity means a sync-looking \code{@isolated(any)} value still uses cross-isolation calling conventions.

\codefile{snippets/isolated-any-await.swift}

\subsection{Example 8. \code{isolated} parameters reuse the ordinary call rules}

An \code{isolated} parameter does not require a new call-rule family. The passed actor witness determines callee isolation, after which the same-isolation and cross-isolation rules apply normally.

\codefile{snippets/isolated-parameter.swift}

\subsection{Example 9. Cross-isolation non-\code{sending} results are rejected}

Cross-isolation non-\code{Sendable} results require explicit \code{sending}. Without that result-side guarantee, derivation fails because the caller cannot safely receive the value.

\codefile{snippets/nonsending-result.swift}

\subsection{Example 10. \code{\#isolation} removes the need for \code{@Sendable}}

When \code{\#isolation} supplies same-isolation evidence to a higher-order helper, non-\code{@Sendable} closure capture and mutable local access remain legal because no cross-isolation boundary is introduced.

\codefile{snippets/isolation-macro.swift}

\subsection{Example 11. \code{async let} temporarily consumes captures with $\undosend$}

\code{async let} creates a child task boundary that temporarily consumes non-\code{Sendable} captures. After \code{await}, captures that were not sent cross-isolation within the child task body are restored via $\undosend$. This distinguishes \code{async let} from \code{Task.init}: the latter either preserves captures through same-actor inheritance or permanently consumes them, whereas \code{async let} supports temporary consumption with restoration.

\subsection{Example 12. Nonisolated sync: ``callable from'' $\neq$ ``convertible to''}

\textsc{call-nonisolated-sync} allows any actor-isolated context to call a nonisolated sync function without \code{await}. However, this call-site property does not yield a new function type conversion edge (\textsc{func-conv}). Closure wrapping \code{\{ g() \}} involves two independent checks: (1)~a \emph{capture check}---whether \code{g}'s region can cross into the closure's isolation---and (2)~a \emph{call check}---whether the closure body can invoke \code{g()}, which always succeeds via \textsc{call-nonisolated-sync}. A \disconnected{} local or a \sending{} parameter passes the capture check; an ordinary parameter at \task{} region does not.

\codefile{snippets/nonisolated-sync-call-vs-capture.swift}

Together, these twelve examples cover the paper's main semantic contrasts: binding versus consumption, same-isolation evidence versus dynamic uncertainty, argument transfer versus result acceptance, temporary versus permanent capture consumption, and call-site compatibility versus type-level convertibility.


\clearpage
\bibliographystyle{ACM-Reference-Format}
\bibliography{references}

\end{document}
